\bookchapter{The ROS 2 input adapter}\label{k:ros}

On the K1, Booster's services publish the robot's state on ROS~2 topics (Chapter~\ref{k:box}). Our program uses ROS~2 for exactly one thing: to \emph{receive} these sensor data. A small \textbf{input adapter} subscribes to the topics in its own thread, keeps the newest values in a buffer, and hands them to the \texttt{SensorInput} module of Chapter~\ref{k:fake}. Nothing else in the code imports \texttt{rclpy}; that is what lets the same code run with the fake robot. Milestone: one.

\begin{needbox}
\begin{itemize}[nosep]
  \item The robot with SSH, your notes on topics and message types (Chapter~\ref{k:box}, Step~6), the joint table (Chapter~\ref{k:model}).
  \item Background: \rtr{ROS 2 Concepts}, \rtr{Writing Your Own Nodes}, \rtr{Logging, Replay and Visualisation}.
\end{itemize}
\end{needbox}

\begin{guessbox}
Topic names (\texttt{/low\_state}, \texttt{/camera/\dots}), message packages (\texttt{booster\_interface}) and field names (\texttt{imu\_state.gyro}, \texttt{motor\_state\_serial[i].q}) below are placeholders\guess. Take the real ones from \verb!ros2 topic info -v! and \verb!ros2 interface show <type>! on your robot. The code in this chapter could not be run without the robot's message packages; it is the structure that matters.
\end{guessbox}

\section*{Design}
\begin{center}
\begin{tikzpicture}[>={Stealth[length=2.2mm]},bx/.style={draw,thick,rounded corners=3pt,minimum width=2.7cm,minimum height=0.85cm,font=\scriptsize,align=center}]
 \node[bx,fill=gray!12] (ros) at (0,0) {ROS 2 topics\\ (Booster services)};
 \node[bx,fill=orange!12] (th) at (3.6,0) {adapter thread\\ \texttt{rclpy} callbacks};
 \node[bx,fill=yellow!20] (buf) at (7.0,0) {\texttt{LatestBuffer}\\ (with a lock)};
 \node[bx,fill=blue!8] (cog) at (10.4,0.7) {cognition loop};
 \node[bx,fill=blue!8] (mot) at (10.4,-0.7) {motion loop};
 \draw[->,thick] (ros)--(th); \draw[->,thick] (th)--node[above,font=\scriptsize]{write}(buf); \draw[->,thick] (buf)--(cog); \draw[->,thick] (buf)--(mot);
\end{tikzpicture}
\captionof{figure}{The ROS boundary. Callbacks only convert and store; they never run robot logic. Both loops read the newest values.}
\end{center}

Rules for the adapter:
\begin{enumerate}[nosep]
  \item \textbf{Convert at the border.} ROS messages are turned into our own dataclasses (with NumPy arrays and SI units) inside the callback. No ROS type leaves the adapter.
  \item \textbf{Reorder at the border.} The joint order of the robot's message may differ from the model's order. Map by name, once, and check it with a test.
  \item \textbf{Keep only the newest, with a time stamp.} The state arrives much faster than cognition runs (for example $500$\,Hz against $20$\,Hz: $25$ messages per cognition cycle). The buffer stores the newest message and the time it arrived; a reader can see how old it is.
  \item \textbf{Quality of service:} \emph{best effort} and \emph{keep last 1} for high-rate sensor data. An old state is worthless, so it must not be queued or re-sent.
\end{enumerate}

\section*{The code}
\begin{verbatim}
import threading, time
import numpy as np
import rclpy
from rclpy.node import Node
from rclpy.qos import QoSProfile, ReliabilityPolicy, HistoryPolicy
from booster_interface.msg import LowState          # placeholder [?]

class LatestBuffer:
    def __init__(self):
        self._lock = threading.Lock()
        self._value, self._t = None, 0.0
    def put(self, value):
        with self._lock:
            self._value, self._t = value, time.monotonic()
    def get(self):
        with self._lock:
            return self._value, self._t

class RosInput(Node):
    def __init__(self, joint_map):
        super().__init__("k1code_input")
        self.joint_map = np.asarray(joint_map)   # model index -> msg index
        self.state = LatestBuffer()
        qos = QoSProfile(reliability=ReliabilityPolicy.BEST_EFFORT,
                         history=HistoryPolicy.KEEP_LAST, depth=1)
        self.create_subscription(LowState, "/low_state", self.on_state, qos)

    def on_state(self, msg):
        motors = msg.motor_state_serial                  # placeholder [?]
        q = np.array([m.q for m in motors])[self.joint_map]
        dq = np.array([m.dq for m in motors])[self.joint_map]
        temp = np.array([m.temperature for m in motors])[self.joint_map]
        self.state.put(SensorFrame(
            t=time.monotonic(), q=q, dq=dq, temperature=temp,
            gyro=np.array(msg.imu_state.gyro),
            gravity=gravity_from_rpy(msg.imu_state.rpy)))

def start_ros_input(joint_map):
    rclpy.init()
    node = RosInput(joint_map)
    threading.Thread(target=rclpy.spin, args=(node,), daemon=True).start()
    return node
\end{verbatim}
\verb!gravity_from_rpy! turns the roll and pitch angles of the IMU into the direction of gravity in the trunk's frame (the yaw does not matter); it is the \emph{projected gravity} that the walking policy uses (\rtr{Velocity tracking and reward design}). Write it with the rotation matrices of \rtr{Rotation and translation matrices}, and test it: upright gives $(0,0,-1)$.

\section*{Steps}
\begin{stepbox}{Step 1: record a bag}
On the robot, while someone moves the hanging robot's limbs by hand (damping): \verb!ros2 bag record -o test_state /low_state! for $20$\,s. Copy it to the laptop. Now you have test data without the robot.
\end{stepbox}
\begin{stepbox}{Step 2: the joint map}
From the message order (documentation or \texttt{ros2 interface show}) and the model's order (\texttt{docs/k1\_joints.txt}), build \texttt{joint\_map} by name. Test: every index appears once; moving the left knee by hand changes the entry called \texttt{left\_knee} in our arrays.
\end{stepbox}
\begin{stepbox}{Step 3: the adapter}
Write \texttt{adapters/ros\_input.py} as above, and an \texttt{RosIO} class that combines it with the SDK output (Chapter~\ref{k:walk}) into the \texttt{RobotIO} interface.
\end{stepbox}
\begin{stepbox}{Step 4: measure}
On the robot, run the program with the real adapter and log, once per second: the number of messages received, and the age of the newest one when cognition reads it.
\end{stepbox}

\begin{checkbox}
\begin{itemize}[nosep]
  \item The message rate matches \verb!ros2 topic hz!. The age of the newest state at reading time is below one period of the topic.
  \item Moving the left knee by hand changes \texttt{left\_knee} and nothing else; the projected gravity is about $(0,0,-1)$ for the hanging robot.
\end{itemize}
\end{checkbox}

\begin{troublebox}
\begin{description}[leftmargin=0.6cm,nosep]
  \item[No messages, but \texttt{ros2 topic echo} works] QoS mismatch: a reliable subscriber does not get data from a best-effort publisher. \verb!ros2 topic info -v! shows the publisher's QoS.
  \item[\texttt{ModuleNotFoundError: booster\_interface}] The Booster message package is not sourced in the environment of the service; the systemd unit needs the ROS environment (an \texttt{EnvironmentFile} or a wrapper script that sources \texttt{setup.bash}).
  \item[The virtual environment cannot see \texttt{rclpy}] \texttt{rclpy} comes from the system's ROS installation, not from PyPI. Create the virtual environment with \verb!--system-site-packages!, and use the same Python version as the robot's ROS.
\end{description}
\end{troublebox}

\begin{keyidea}[What to remember]
\begin{itemize}[nosep]
  \item ROS~2 only at the input border, in its own thread; convert and reorder in the callback; store the newest value with a time stamp.
  \item Best effort, keep last 1 for high-rate sensor data.
  \item A recorded bag is test data for the adapter.
\end{itemize}
\end{keyidea}
